% ===================== SLIDE 16 — CONTRIBUTION 1 =====================
\begin{frame}[t]
\pmheader{CONTRIBUTION 1}{18 / 22}
\par\pmtitle{How do information costs shape participation?}
\begin{pmbody}
{\color{pmGrey}\fontsize{20}{27}\selectfont How do changes in information costs affect who acquires information and who trades?\par}
\begin{itemize}\pmspreaditems
  \item Lower costs may allow additional participants to become informed.
  \item Changes in relative information advantages may alter entry, exit, and trading intensity.
  \item The resulting participation changes may affect both information production and liquidity supply.
  \item AI is one possible setting. Other candidates include changes in data access or information disclosure, provided they offer credible identification.
\end{itemize}
\begin{pmcallout}[Evidence needed]
  A measurable change in information costs and a design that separates its effects from concurrent changes.
\end{pmcallout}
{\color{pmGrey}\fontsize{11}{14}\selectfont References: Theoretical motivation: Grossman and Stiglitz (1980).\par}
\end{pmbody}
\end{frame}

% ===================== SLIDE 19 — CONTRIBUTION 2 =====================
\begin{frame}[t]
\pmheader{CONTRIBUTION 2}{19 / 22}
\par\pmtitle{Which trading functions improve accuracy?}
\begin{pmbody}
{\color{pmGrey}\fontsize{20}{27}\selectfont Which trading functions account for the relationship between skilled participation and accurate prices?\par}
\begin{itemize}\pmspreaditems
  \item Skilled traders may contribute through aggressive trades that incorporate information into prices.
  \item They may also supply liquidity that allows other participants to trade.
  \item Execution advantages, inventory risk, and information advantages may explain different components of profitability.
  \item Persistent liquidity provision may matter differently from occasionally submitting a limit order.
\end{itemize}
\begin{pmcallout}[Evidence needed]
  Measures that distinguish forecasting skill, execution role, and sustained liquidity provision, ideally combined with independent variation in participation.
\end{pmcallout}
{\color{pmGrey}\fontsize{11}{14}\selectfont References: Akey et al. (2026); Heesen and Luong (2026).\par}
\end{pmbody}
\end{frame}

% ===================== SLIDE 18 — CONTRIBUTION 3 =====================
\begin{frame}[t]
\pmheader{CONTRIBUTION 3}{20 / 22}
\par\pmtitle{When does price consensus signal better information?}
\begin{pmbody}
{\color{pmGrey}\fontsize{20}{27}\selectfont More similar prices can arise from different underlying mechanisms.\par}
\begin{itemize}\pmspreaditems
  \item Trading may successfully incorporate dispersed information.
  \item Participants may instead rely on common signals and share similar errors.
  \item Changes in the composition or timing of trades may also reduce measured price dispersion.
\end{itemize}
\begin{pmcallout}[Potential contribution]
  Distinguish these mechanisms and establish when stronger price consensus predicts better calibration.
\end{pmcallout}
\begin{pmcallout}[Evidence needed]
  Transaction data linked to realized outcomes, with comparable prediction horizons and attention to changes in trading composition.
\end{pmcallout}
{\color{pmGrey}\fontsize{11}{14}\selectfont References: Ostrovsky (2012); Page and Clemen (2013).\par}
\end{pmbody}
\end{frame}
